\documentclass[conference]{IEEEtran}
\IEEEoverridecommandlockouts
\usepackage{cite}
\usepackage{amsmath,amssymb,amsfonts}
\usepackage{graphicx}
\usepackage{textcomp}
\usepackage{xcolor}
\usepackage{booktabs}
\usepackage{siunitx}
\def\BibTeX{{\rm B\kern-.05em{\sc i\kern-.025em b}\kern-.08em
    T\kern-.1667em\lower.7ex\hbox{E}\kern-.125emX}}
\begin{document}

\title{EEE4120F Practical 2: Introduction to MATLAB Parallel
Computing Toolkit (PCT)}

\author{\IEEEauthorblockN{Matteo Buxman}
\IEEEauthorblockA{BXMMAT001}
\and
\IEEEauthorblockN{Emmanuel Basua}
\IEEEauthorblockA{BSXEMM001}
}

\maketitle

\begin{abstract}
This report analyses the performance of serial and parallel Mandelbrot set
generation in MATLAB using the Parallel Computing Toolbox (PCT).  Eight
resolutions (SVGA through 8K~UHD) were benchmarked with a fixed iteration
limit of 1\,000 and worker pools of 2--6.  Speedup, efficiency, and
Amdahl's Law serial-fraction estimates are reported.  Results show that
meaningful speedup requires sufficiently large problem sizes, with best
observed speedup of $2.46\times$ at 8K, while small resolutions suffer from
parallelisation overhead.  Empirically fitted serial fractions decrease from
60\% (SVGA) to 37--38\% (5K/8K), consistent with Gustafson's Law.
\end{abstract}

\section{Introduction}
Modern multi-core processors offer significant performance potential when
software exploits parallelism.  MATLAB's Parallel Computing Toolbox (PCT)
provides the \texttt{parfor} construct to distribute loop iterations across
a pool of workers without explicit thread management.

The Mandelbrot set is an ideal benchmark for evaluating parallel performance:
each pixel is computed independently, making the problem embarrassingly
parallel.  However, per-pixel computational cost varies widely---pixels near
the set boundary iterate up to the maximum count, while exterior pixels escape
in very few iterations---creating a natural load-balancing challenge that
tests how well the parallelisation framework handles heterogeneous workloads.

Formally, the set is defined by iterating $f_c(z) = z^2 + c$ from $z = 0$
for each point $c \in \mathbb{C}$ in the complex plane.  A point belongs to
the set if the orbit $\{z_n\}$ remains bounded; in practice, a pixel is
coloured by the number of iterations before $|z| > 2$, capped at
\texttt{max\_iterations} = 1\,000.

This practical aims to (i)~implement serial and parallel Mandelbrot
generators in MATLAB, (ii)~benchmark both across eight standard resolutions
from SVGA to 8K~UHD, (iii)~quantify speedup and efficiency as functions of
resolution and worker count, and (iv)~estimate the parallelisable fraction
of the workload via Amdahl's Law and observe how this fraction changes with
problem size.

\section{Methodology}

\subsection{Algorithm}

The complex plane is sampled over the standard Mandelbrot viewing region:
real axis $x \in [-2.0,\,0.5]$, imaginary axis $y \in [-1.2,\,1.2]$.
Each pixel at grid position $(col, row)$ maps to a complex constant
$c = x_0 + iy_0$.  The escape-time algorithm iterates
\begin{equation}
  z_{n+1} = z_n^2 + c, \quad z_0 = 0
\end{equation}
until $|z_n|^2 > 4$ or the iteration count reaches \texttt{max\_iterations}
= 1\,000.  The resulting iteration grid is colour-mapped using MATLAB's
\texttt{hot} colormap and saved as a PNG via \texttt{imwrite}.

\subsection{Serial Implementation}

\texttt{mandelbrot\_sequential.m} computes the Mandelbrot set using two
nested \texttt{for} loops: the outer loop iterates over image rows
(imaginary axis) and the inner loop over columns (real axis).  The
computation is isolated in a helper function \texttt{compute\_mandelbrot},
separating the numerical kernel from the I/O layer.  This design provides a
clean, deterministic baseline: no worker overhead, no inter-process
communication, and fully deterministic execution order, ensuring that all
parallel speedup measurements are relative to a well-characterised reference.

\subsection{Parallel Implementation}

\texttt{mandelbrot\_parallel.m} replaces the outer row loop with
\texttt{parfor}, allowing each worker to independently compute one or more
rows of the iteration grid.  The inner column loop remains a standard
\texttt{for} loop, as MATLAB does not support nested \texttt{parfor}.

Correct \texttt{parfor} variable classification ensures both correctness and
performance: \texttt{iteration\_grid(row,:)} is a sliced output (one writer
per iteration), \texttt{y\_range(row)} a sliced input, \texttt{x\_range} and
\texttt{max\_iterations} are broadcast once to all workers, and loop
temporaries (\texttt{x0}, \texttt{y0}, etc.) are worker-private.  A local
accumulator \texttt{row\_result} buffers inner-loop output before the sliced
assignment, improving cache locality per worker.

Worker pool management is handled in \texttt{run\_analysis.m}: a
\texttt{parpool} is created (or resized) \emph{before} timing begins,
ensuring that pool spin-up latency---typically 5--15\,s---is excluded from
all reported execution times.  Pools are reused across resolutions within
the same worker-count sweep to reflect realistic sustained usage.

\subsection{Timing and Metrics}

All timings use MATLAB's \texttt{tic}/\texttt{toc}, wrapping only the
compute-and-save function call.  Serial runs precede all parallel runs so
that their wall-clock times serve as the reference denominator for speedup.
Speedup and efficiency are defined as:
\begin{equation}
  S = \frac{T_{\text{serial}}}{T_{\text{parallel}}}, \qquad
  E = \frac{S}{P} \times 100\%
\end{equation}

Amdahl's Law gives the theoretical maximum speedup with $P$ workers when a
fraction $f$ of the workload is parallelisable:
\begin{equation}
  S_{\max} = \frac{1}{(1 - f) + \dfrac{f}{P}}
  \label{eq:amdahl_s}
\end{equation}
Rearranging to solve for $f$ given measured speedup $S$ at $P$ workers:
\begin{equation}
  f = \frac{1 - \tfrac{1}{S}}{1 - \tfrac{1}{P}}
  \label{eq:amdahl_f}
\end{equation}
For each resolution, $f$ is computed at every worker count
$P \in \{2,\ldots,6\}$ and averaged, yielding an empirical per-resolution
estimate that accounts for measurement variability across different degrees
of parallelism.

\section{Results}

\subsection{Baseline Execution Time}

Table~\ref{tab:serial} presents the serial baseline execution times across
all eight resolutions.  These serve as the reference denominator for all
parallel performance calculations.

\begin{table}[h]
\centering
\caption{Serial (baseline) execution times}
\label{tab:serial}
\renewcommand{\arraystretch}{1.15}
\begin{tabular}{@{}lrr@{}}
\toprule
\textbf{Resolution} & \textbf{Pixels (MP)} & \textbf{Time (s)} \\
\midrule
SVGA (800$\times$600)      & 0.48  & 0.67  \\
HD (1280$\times$720)       & 0.92  & 0.83  \\
FHD (1920$\times$1080)     & 2.07  & 1.86  \\
2K (2048$\times$1080)      & 2.21  & 1.93  \\
QHD (2560$\times$1440)     & 3.69  & 3.30  \\
4K UHD (3840$\times$2160)  & 8.29  & 7.32  \\
5K (5120$\times$2880)      & 14.75 & 13.20 \\
8K UHD (7680$\times$4320)  & 33.18 & 30.10 \\
\bottomrule
\end{tabular}
\end{table}

Execution time scales approximately linearly with pixel count, confirming
$O(N)$ complexity.  The 8K image (33.18\,MP) takes 30.10\,s, roughly
$45\times$ longer than SVGA (0.48\,MP, 0.67\,s).

\subsection{Speedup}

Figure~\ref{fig:speedup} shows measured speedup versus worker count for all
eight resolutions, with Amdahl's Law theoretical curves (dashed) overlaid.
The theoretical curves are generated using~(\ref{eq:amdahl_s}) with
per-resolution parallel fractions empirically fitted via~(\ref{eq:amdahl_f}).

\begin{figure}[h]
\centering
\includegraphics[width=\linewidth]{figures/speedup_vs_workers.png}
\caption{Measured speedup vs.\ worker count with Amdahl's Law theoretical
predictions (dashed) using empirically fitted parallel fractions per
resolution.  The white dashed line represents ideal linear scaling
($S = P$).}
\label{fig:speedup}
\end{figure}

All resolutions fall well below the ideal $S = P$ line.  The best observed
speedup is $2.46\times$ at 8K with 6 workers, against a theoretical ideal
of $6\times$.  Larger resolutions consistently achieve higher speedup as
the computation-to-overhead ratio improves: at 6 workers, speedup increases
monotonically from $1.56\times$ (SVGA) to $2.46\times$ (8K).

SVGA is the most striking outlier---it peaks at $1.61\times$ (at 4 workers)
and shows performance degradation beyond that point.  At 6 workers, SVGA
achieves only $1.56\times$, meaning the additional two workers contribute
negative marginal speedup.  This confirms that at sub-second workloads
(0.67\,s serial), \texttt{parfor} scheduling cost dominates.

The Amdahl's Law curves fit the measured data closely across most
resolutions, validating the empirical parallel-fraction estimation approach.
Deviations at low worker counts for small resolutions (e.g.\ SVGA, HD) arise
from measurement noise at short runtimes, where OS scheduling jitter
represents a larger fraction of total execution time.

\subsection{Parallel Efficiency}

Figure~\ref{fig:efficiency} shows parallel efficiency across resolutions
for each worker count.

\begin{figure}[h]
\centering
\includegraphics[width=\linewidth]{figures/efficiency_vs_resolution.png}
\caption{Parallel efficiency vs.\ resolution for 2--6 workers.  The dashed
line at 100\% represents ideal scaling.}
\label{fig:efficiency}
\end{figure}

Efficiency ranges from approximately 26\% (SVGA at 6 workers) to 67\%
(HD at 2 workers), consistently below the 100\% ideal.  Two clear trends
are visible.  First, efficiency generally improves with resolution as the
per-worker workload grows relative to the fixed parallelisation overhead,
confirming that larger problem sizes amortise communication costs more
effectively.  Second, efficiency decreases with worker count---a direct
consequence of Amdahl's Law, as adding workers dilutes the parallel portion
while the serial component remains constant.

The 4K~UHD resolution achieves the best efficiency at higher worker counts
(46.0\% at 5 workers), suggesting its problem size strikes the best balance
between amortising \texttt{parfor} overhead and avoiding memory bandwidth
saturation.  At 8K, despite having the highest speedup, efficiency at 6
workers (41.0\%) is lower than 4K's, hinting that memory contention between
workers begins to limit scaling at very large grid sizes where the iteration
matrix exceeds 250\,MB.

\subsection{Amdahl's Law --- Parallel Fraction}

Figure~\ref{fig:par_frac} shows the empirically estimated parallel fraction
$f$ for each resolution, computed by averaging~(\ref{eq:amdahl_f}) across
all five worker counts $P \in \{2,\ldots,6\}$.

\begin{figure}[h]
\centering
\includegraphics[width=\linewidth]{figures/amdahls_law_serial_fraction.png}
\caption{Estimated serial fraction $(1-f)$ per resolution, fitted from
measured speedup data using Amdahl's Law~(\ref{eq:amdahl_f}).  Lower values
indicate a larger parallelisable share of the workload.}
\label{fig:par_frac}
\end{figure}

The serial fraction $(1-f)$ is highest at SVGA (60.1\%), where
parallelisation overhead dominates, then drops sharply to the 40--44\%
range for mid-range resolutions (HD through QHD), before settling at
36.8\% (5K) and 38.0\% (8K).  The overall trend from 60\% down to
${\sim}37\text{--}38\%$ is significant: under a strict interpretation of
Amdahl's Law, $f$ should be a fixed property of the algorithm, independent
of problem size.  The observed decrease---particularly the large gap between
SVGA and the remaining resolutions---reveals that what appears as ``serial''
overhead is dominated by \emph{fixed-cost} components (\texttt{parfor}
dispatch, inter-process data transfer, and the serialised \texttt{imwrite}
call) rather than an inherently sequential portion of the Mandelbrot
computation itself.

As resolution increases, these fixed costs become a smaller fraction of
total runtime, effectively increasing the parallelisable share.  This
behaviour is consistent with \textbf{Gustafson's Law}, which predicts that
the parallel fraction of a workload grows with problem size because the
serial component remains constant while total work scales.  Gustafson's
scaled speedup is defined as:
\begin{equation}
  S_G = P + (1 - P) \cdot (1 - f)
  \label{eq:gustafson}
\end{equation}
where $(1-f)$ is the serial fraction at a given problem size.  Unlike
Amdahl's Law, which assumes a fixed total workload and predicts diminishing
returns from adding processors, Gustafson's model captures the practical
reality that users typically scale problem size with available resources,
making parallelism increasingly effective for larger workloads.

\section{Discussion}

\textbf{Sub-linear speedup and the serial bottleneck.}
Across all resolutions, measured speedup peaks at $2.46\times$ with 6
workers, well below the theoretical ideal of $6\times$.  This implies that
a significant serial component persists throughout execution.  The primary
contributors are: (i)~the serialised image write in
\texttt{mandelbrot\_plot}, which cannot overlap with parallel computation;
(ii)~\texttt{parfor} row-distribution and result-collection overhead, where
MATLAB must coordinate work assignment and gather partial results from each
worker; and (iii)~inter-process memory copies of broadcast variables,
particularly \texttt{x\_range}.  At 8K, the estimated serial fraction of
38.0\% sets a theoretical Amdahl ceiling of ${\sim}2.6\times$ even with
infinitely many workers---a hard limit that cannot be overcome without
restructuring the serial components themselves.

\textbf{Overhead dominance at small resolutions.}
SVGA is the most instructive case: its best speedup of $1.61\times$ at 4
workers demonstrates that at sub-second workloads, \texttt{parfor} fixed
costs (worker dispatch, row slicing, result aggregation) consume a
disproportionate fraction of total runtime.  Beyond 4 workers, each
additional worker adds more coordination overhead than it saves in
computation, yielding \emph{negative} marginal speedup.  This underscores
a fundamental principle of parallel computing: parallelisation is only
beneficial when the granularity of the parallelisable work is large enough
to amortise the fixed cost of distributing and collecting that work.

\textbf{Scalability with problem size.}
The decreasing serial fraction with resolution
(Fig.~\ref{fig:par_frac})---from 60.1\% at SVGA to ${\sim}37\text{--}38\%$
at 5K and 8K---demonstrates that Amdahl's Law, while accurate for predicting
speedup at a fixed problem size, provides an overly pessimistic view of
parallel scalability when problem size is allowed to grow.  The data aligns
with Gustafson's Law: increasing resolution grows the parallel workload
while leaving the fixed overhead largely unchanged, yielding progressively
better speedups and higher effective parallel fractions.  The fact that the
serial fraction plateaus around 37--38\% for the largest resolutions
suggests that an irreducible serial component (likely the serialised image
write) persists even after \texttt{parfor} overhead is fully amortised.

\textbf{Load imbalance and scheduling.}
The Mandelbrot set has inherently unequal per-pixel cost: pixels near the
set boundary iterate to \texttt{max\_iterations} = 1\,000, while exterior
pixels escape in fewer than 10 iterations.  To quantify this,
Figure~\ref{fig:load_per_row} profiles per-row computation time at 4K~UHD
with 6 workers.  The characteristic ``mountain'' shape confirms that rows
intersecting the main cardioid and period-2 bulb (roughly rows 400--1500)
are orders of magnitude more expensive than edge rows.

\begin{figure}[h]
\centering
\includegraphics[width=\linewidth]{figures/load_balance_per_row.png}
\caption{Per-row computation time at 4K~UHD coloured by assigned worker.
Rows near the Mandelbrot boundary (400--1500) are significantly more
expensive.  Contiguous colour bands reveal that \texttt{parfor} dispatches
iterations in fixed-size chunks rather than individual rows.}
\label{fig:load_per_row}
\end{figure}

The colour bands in Figure~\ref{fig:load_per_row} reveal that MATLAB's
\texttt{parfor} scheduler divides the iteration space into fixed-size
contiguous chunks and dispatches them to workers as they become
available~\cite{parfor_schedule}.  Workers that receive chunks of cheap edge
rows finish quickly and request more chunks, accumulating a high row count
but low total computation time.  Workers assigned chunks in the expensive
boundary region remain busy for longer, completing fewer rows but
accumulating far more computation time.

Figure~\ref{fig:load_total} quantifies the resulting imbalance: the busiest
worker (W3, 275 rows, 1.83\,s) accumulated ${\sim}2.8\times$ more
computation time than the lightest (W2, 678 rows, 0.65\,s), despite W2
processing over twice as many rows.  This confirms that chunk-level
granularity is too coarse to balance the load when per-row costs vary as
dramatically as they do in the Mandelbrot set.  Because overall wall-clock
time is bounded by the \emph{slowest} worker, this imbalance directly
limits achievable speedup.  Finer-grained scheduling---dispatching
individual rows or small tiles rather than large contiguous
chunks---would improve worker utilisation.

\begin{figure}[h]
\centering
\includegraphics[width=\linewidth]{figures/load_balance_total.png}
\caption{Total computation time per worker at 4K~UHD with 6 workers.  Row
counts are labelled above each bar.  Workers assigned expensive boundary
rows (W3) accumulate far more computation time than workers processing
cheap edge rows (W2), despite handling fewer rows.}
\label{fig:load_total}
\end{figure}

\textbf{Limitations.}
All results are single-run measurements without repeated trials or
statistical averaging.  Transient OS activity, MATLAB JIT warm-up effects,
and memory allocation variability introduce noise that is not filtered out.
Additionally, the experiment was limited to 6 workers by available hardware;
testing with higher worker counts would better reveal whether the measured
speedup converges to the Amdahl ceiling or whether additional scaling
is achievable.

\section{Conclusion}

Three key findings emerge from this analysis:

\begin{enumerate}
  \item \textbf{Parallelism is resolution-dependent.}  Meaningful speedup
        ($>2\times$) requires at least Full~HD resolution (2.07\,MP).  At
        SVGA (0.48\,MP), \texttt{parfor} overhead exceeds the computational
        benefit, and adding workers beyond 4 actively degrades performance.

  \item \textbf{Speedup is bounded by the serial fraction.}  The best
        observed speedup of $2.46\times$ at 8K with 6 workers corresponds
        to a serial fraction of 38.0\% under Amdahl's Law, setting a
        hard theoretical ceiling of ${\sim}2.6\times$ regardless of worker
        count.  The serial bottleneck is primarily the serialised image
        write and \texttt{parfor} coordination overhead.

  \item \textbf{Gustafson's Law applies at scale.}  The empirically fitted
        serial fraction decreases from 60.1\% (SVGA) to ${\sim}37\text{--}38\%$
        (5K/8K), demonstrating that fixed overheads become a smaller share
        of total work as problem size grows.  The plateau at
        ${\sim}37\text{--}38\%$ for the largest resolutions indicates an
        irreducible serial component, but parallelism nonetheless becomes
        increasingly effective as resolution scales.
\end{enumerate}

Potential improvements include replacing \texttt{imwrite} with asynchronous
I/O via \texttt{parfeval} to overlap file writes with computation,
vectorising the inner column loop using MATLAB array operations to reduce
per-pixel loop overhead, and implementing tile-based work distribution for
better load balancing across the Mandelbrot set boundary.

\begin{thebibliography}{9}
\bibitem{pct_tutorial}
S.~Winberg and T.~Webb, ``MATLAB Parallel Computing Toolkit (PCT): A Practical
Guide for Practical 2,'' EEE4120F High Performance Embedded Systems Reference
Guide, University of Cape Town, 2026.

\bibitem{pct_manual}
S.~Winberg and T.~Webb, ``Practical 2: Introduction to MATLAB Parallel
Computing Toolkit (PCT),'' EEE4120F High Performance Embedded Systems
Practical Manual, University of Cape Town, 2026.

\bibitem{parfor_schedule}
E.~Flagg, ``parfor schedule not very smart,'' MATLAB Answers, MathWorks,
2013. [Online]. Available:
\texttt{https://www.mathworks.com/matlabcentral/answers/85783}
\end{thebibliography}

\end{document}
